\documentclass[a4paper, 12pt, twoside]{article}
\usepackage{ifxetex}
\ifxetex
  \usepackage{fontspec}
\else
  \usepackage[T1]{fontenc}
  \usepackage[utf8]{inputenc}
\fi
\usepackage{babel}
\usepackage{lmodern}
\usepackage{graphicx}
\usepackage{enumitem}
\usepackage{lscape}
\usepackage{subcaption}
\usepackage{imakeidx}
\usepackage{tocbibind}
\usepackage[hidelinks]{hyperref}

\usepackage{xcolor}
\usepackage{listings}
\lstset{%
	basicstyle=\footnotesize\ttfamily\linespread{1}\selectfont,%
	numbers=left,%
	backgroundcolor=\color{lightgray},%
	showstringspaces=false,%
	breaklines=true%
}
\usepackage[newfloat]{minted}
\SetupFloatingEnvironment{listing}{}
\usemintedstyle{emacs}
\setminted{linenos, breaklines, tabsize=4, bgcolor=lightgray}

\usepackage{csquotes}
\usepackage[normalem]{ulem}

\usepackage[margin=2.5cm]{geometry}
\usepackage{setspace}
\setlength\parindent{1cm}
\onehalfspacing

\begin{document}

\begin{titlepage}
    \vspace*{5cm}
    \begin{center}
        \Huge \textbf{CLI and TUI Parameters (Step by Step)}
    \end{center}
    \vfill
\end{titlepage}

\tableofcontents
\clearpage

\subsection*{About}
\addcontentsline{toc}{subsection}{About}

This guide explains the command:

\begin{listing}[h!]
    \begin{minted}{bash}
md2tex source.md -c -t ./utils/template.tex -o ./export/export.tex -d book -u -f
    \end{minted}
\end{listing}

and how to set the same options in the TUI.

\subsection*{1. Base input}
\addcontentsline{toc}{subsection}{1. Base input}

\begin{listing}[h!]
    \begin{minted}{bash}
md2tex source.md
    \end{minted}
\end{listing}

\begin{itemize}
    \item \texttt{source.md} is the required input Markdown file.
    \item It must exist and use \texttt{.md} extension.
\end{itemize}

TUI equivalent:

\begin{itemize}
    \item \texttt{Input Markdown File (*.md) [required]} = \texttt{source.md}
\end{itemize}

\subsection*{2. Build a complete TeX document}
\addcontentsline{toc}{subsection}{2. Build a complete TeX document}

\begin{listing}[h!]
    \begin{minted}{bash}
-c
    \end{minted}
\end{listing}

\begin{itemize}
    \item Enables complete TeX output (preamble + \texttt{\textbackslash\{\}\textbackslash\{\}begin\{document\}} + \texttt{\textbackslash\{\}\textbackslash\{\}end\{document\}} + TOC support).
    \item Without \texttt{-c}, output is body-only TeX content.
\end{itemize}

TUI equivalent:

\begin{itemize}
    \item Enable \texttt{Complete TeX file (-c)}
\end{itemize}

\subsection*{3. Use a custom TeX template}
\addcontentsline{toc}{subsection}{3. Use a custom TeX template}

\begin{listing}[h!]
    \begin{minted}{bash}
-t ./utils/template.tex
    \end{minted}
\end{listing}

\begin{itemize}
    \item Uses a custom template instead of the bundled default template.
    \item The template must include \texttt{@@BODYTOKEN@@}.
    \item Common pattern: keep preamble/style in template and inject converted Markdown body at \texttt{@@BODYTOKEN@@}.
\end{itemize}

TUI equivalent:

\begin{itemize}
    \item \texttt{Template File (*.tex) [optional]} = \texttt{./utils/template.tex}
\end{itemize}

\subsubsection*{What is a template in md2tex?}
\addcontentsline{toc}{subsubsection}{What is a template in md2tex?}

A template is a complete \texttt{.tex} skeleton that \texttt{md2tex} uses as a wrapper around converted Markdown content.

Think of it as:

\begin{itemize}
    \item Your LaTeX layout/style (preamble, fonts, margins, packages)
    \item Plus one placeholder where converted Markdown will be inserted
\end{itemize}

Required token:

\begin{itemize}
    \item \texttt{@@BODYTOKEN@@} (mandatory): replaced with converted body content
\end{itemize}

Optional tokens supported by md2tex default flow:

\begin{itemize}
    \item \texttt{@@DOCUMENTCLASSTOKEN@@}: replaced with \texttt{article} or \texttt{book}
    \item \texttt{@@TITLETOKEN@@}: replaced with the generated title page (if first Markdown heading is level 1)
    \item \texttt{@@CODEPACKAGES@@}: replaced with code package setup (\texttt{listings}/\texttt{minted} + \texttt{xcolor}) according to backend strategy
\end{itemize}

Minimal valid custom template example:

\begin{listing}[h!]
    \begin{minted}{tex}
\\documentclass{@@DOCUMENTCLASSTOKEN@@}
@@CODEPACKAGES@@
\\begin{document}
@@TITLETOKEN@@
@@BODYTOKEN@@
\\end{document}
    \end{minted}
\end{listing}

If \texttt{@@BODYTOKEN@@} is missing, conversion fails because md2tex has nowhere to inject the Markdown output.

\subsection*{4. Choose custom output path}
\addcontentsline{toc}{subsection}{4. Choose custom output path}

\begin{listing}[h!]
    \begin{minted}{bash}
-o ./export/export.tex
    \end{minted}
\end{listing}

\begin{itemize}
    \item Writes output to \texttt{./export/export.tex}.
    \item If directory does not exist, it is created.
    \item If extension is missing, \texttt{.tex} is appended.
\end{itemize}

TUI equivalent:

\begin{itemize}
    \item \texttt{Output TeX File (*.tex) [optional]} = \texttt{./export/export.tex}
\end{itemize}

\subsection*{5. Select document class}
\addcontentsline{toc}{subsection}{5. Select document class}

\begin{listing}[h!]
    \begin{minted}{bash}
-d book
    \end{minted}
\end{listing}

\begin{itemize}
    \item Sets TeX document class to \texttt{book}.
    \item Allowed values: \texttt{article} or \texttt{book}.
    \item Affects heading mapping (\texttt{\#} as \texttt{\textbackslash\{\}\textbackslash\{\}chapter\{\}} in \texttt{book}, \texttt{\textbackslash\{\}\textbackslash\{\}section\{\}} in \texttt{article}).
\end{itemize}

TUI equivalent:

\begin{itemize}
    \item \texttt{document\_class} selector = \texttt{book}
\end{itemize}

\subsection*{6. Convert headings as unnumbered}
\addcontentsline{toc}{subsection}{6. Convert headings as unnumbered}

\begin{listing}[h!]
    \begin{minted}{bash}
-u
    \end{minted}
\end{listing}

\begin{itemize}
    \item Uses unnumbered section commands (e.g., \texttt{\textbackslash\{\}\textbackslash\{\}chapter\textit{\{\}}, \texttt{\textbackslash\{\}\textbackslash\{\}section}\{\}}).
    \item The converter still adds TOC entries with \texttt{\textbackslash\{\}\textbackslash\{\}addcontentsline}.
\end{itemize}

TUI equivalent:

\begin{itemize}
    \item Enable \texttt{Unnumbered headers (-u)}
\end{itemize}

\subsection*{7. Convert inline quotes to French quotes}
\addcontentsline{toc}{subsection}{7. Convert inline quotes to French quotes}

\begin{listing}[h!]
    \begin{minted}{bash}
-f
    \end{minted}
\end{listing}

\begin{itemize}
    \item Converts inline quoted text to \texttt{\textbackslash\{\}\textbackslash\{\}enquote\{...\}} (via \texttt{csquotes}).
\end{itemize}

TUI equivalent:

\begin{itemize}
    \item Enable \texttt{French quotes (-f)}
\end{itemize}

\subsection*{Full mapping summary}
\addcontentsline{toc}{subsection}{Full mapping summary}

\begin{itemize}
    \item CLI \texttt{source.md} -> TUI \texttt{Input Markdown File}
    \item CLI \texttt{-c} -> TUI \texttt{Complete TeX file (-c)}
    \item CLI \texttt{-t PATH} -> TUI \texttt{Template File}
    \item CLI \texttt{-o PATH} -> TUI \texttt{Output TeX File}
    \item CLI \texttt{-d book|article} -> TUI \texttt{document\_class}
    \item CLI \texttt{-u} -> TUI \texttt{Unnumbered headers (-u)}
    \item CLI \texttt{-f} -> TUI \texttt{French quotes (-f)}
\end{itemize}

\subsection*{Extra options you can also use}
\addcontentsline{toc}{subsection}{Extra options you can also use}

CLI:

\begin{listing}[h!]
    \begin{minted}{bash}
--code-backend auto|minted|listings
--shell-escape
    \end{minted}
\end{listing}

TUI equivalents:

\begin{itemize}
    \item \texttt{Enable shell-escape} toggle
    \item Code backend is \texttt{auto} by default:
    \begin{itemize}
        \item shell-escape enabled -> \texttt{minted}
        \item shell-escape disabled -> \texttt{listings}
    \end{itemize}
\end{itemize}

\subsection*{Practical examples}
\addcontentsline{toc}{subsection}{Practical examples}

\begin{enumerate}
    \item Minimal conversion:
\end{enumerate}

\begin{listing}[h!]
    \begin{minted}{bash}
md2tex myfiles/example.md
    \end{minted}
\end{listing}

\begin{enumerate}
    \item Complete book with custom template and output:
\end{enumerate}

\begin{listing}[h!]
    \begin{minted}{bash}
md2tex myfiles/example.md -c -t ./templates/book.tex -o ./build/notes.tex -d book -u -f
    \end{minted}
\end{listing}

\begin{enumerate}
    \item Force minted with shell-escape expectations:
\end{enumerate}

\begin{listing}[h!]
    \begin{minted}{bash}
md2tex myfiles/example.md -c --code-backend minted --shell-escape
    \end{minted}
\end{listing}


\end{document}
